\section{作品目标与设计思路}
\subsection{作品定位}
嘟嘟哒 是面向持续群体交流的社交智能体。其产品定位为“生活在 QQ 群里的 AI 群友”，技术目标是在同一条运行链中完成理解群聊情境、判断参与方式、使用外部能力、形成回答并可靠投递。校园生活与教学支持是首批应用场景，系统的接口划分同时支持向其他群体服务扩展。

群聊具有与单人问答不同的条件：当前消息可能只是一句追问，指向另一人的发言；同一时间可能出现多个话题；群成员之间的更正具有先后关系；一次有用的回答也可能在错误的时机造成打扰。因此，作品把“是否回应、回应谁、需要什么资料、如何表达”分别建模，使技术实现能够围绕真实交互问题调整。

\subsection{需要解决的问题}
\begin{enumerate}
\item \textbf{信息分散。}课程评论、活动、教务安排和通知分布于不同系统，用户需要反复切换页面，并自行核对课程身份、时间范围和资料版本。
\item \textbf{上下文不完整。}省略、引用、别称和更正使当前短句难以独立理解；仅拼接历史又可能把旧安排当成现行事实。
\item \textbf{参与缺乏节奏。}逐条回复会干扰讨论，单纯等待指令则难以利用连续讨论中已经形成的需求。
\item \textbf{难以解释运行过程。}模型、查询服务、配置和消息平台分别出错时，管理员需要知道具体失败环节，以及保存的配置是否实际生效。
\end{enumerate}

\subsection{设计原则与工程选择}
\textbf{明确分工。}输入适配器确认平台事件，感知模块解释语义，社交策略选择行动，能力层获取资料，回答层组织内容。身份、群范围、执行权限和输出目标由程序维护，模型承担语义提取与文本生成。

\textbf{以业务能力连接异构服务。}对智能体暴露的是课程查询、活动搜索、通知读取等业务单元。具体 MCP 工具名、参数投影和服务会话留在适配层。这样既能复用校验与错误处理，也便于保留各数据源的差异。

\textbf{按任务分配计算。}模型档位、推理强度与输出篇幅独立配置，避免将“详细回答”直接解释为“必须换更大的模型”。感知与回答分别选择模型角色和预算。

\textbf{让运行可观察。}QQ 与 Web 预览共用 Runtime；控制台展示实际运行模型、能力、耗时及结果。开发验证使用同一套编排和可替换适配器，减少独立演示实现带来的偏差。

\subsection{本次报告的实现口径}
报告以 2026 年 9 月 6 日本地主仓为对象。开始核对时主仓提交为 \code{2b52593}，工作区还包含录制与查询相关的在途改动；后续章节同时引用明确日期的已存验证材料。提交号用于定位源码，不表示所有历史测试都运行于该提交，也不表示本次重新发布了程序。

本文将实现分为三类：\textbf{已进入当前运行组合}，如校园查询、模型路由和按群自动参与；\textbf{已完成独立组件或离线验证}，如 Memory v2 和 Bandit 估计器；\textbf{后续设计}，如通用多步 Skill、长期群友印象与自动学习。各项验证的运行模式在第\ref{sec:evaluation}节单独列出。
\source{[S1, S2, S3, S10]。历史原型、研究报告与当前主仓具有不同适用范围。}
